% Notes of the Montreal seminar, 1970 (cote 129, typed part): archive cote 129, site page 52 (PDF page 53).
% typescript. Draft compared: none.
% AI-transcribed by Claude Opus 5.5 (claude-opus-5-5, adaptive thinking, effort high), Message Batches API, batch mtl129_full (msgbatch_01AGnDKHwy6ATCM7ehUa6263), 2026-11-13; not yet checked by eye.
\scanpage{129/52}
\mspage{\hand{II 24}}

Soit $\{x_1,\dots,x_r\}$ un système de générateurs du W-module $M$ ;
les relations vérifiées dans l'anneau $A$ entraînent que l'on peut trou-%
ver un système de monômes en les variables $T_{x_1},\dots,T_{x_r}$ , $T_{Vx_1},\dots,T_{Vx_r},\dots,$
$T_{V^Nx_1},\dots,T_{V^Nx_r}$ et de degrés bornés, tels que cette famille de monômes
engendre $A$ en tant que k-espace vectoriel.

D'après 2.5.2 il suffit de montrer que $\mathbb{V}_{\mathbb{E}(M)} = \mathbb{V}$ est nilpo-%
tent pour prouver que $\mathbb{E}(M)$ est unipotent. Comme $\mathbb{E}(M)$ est un sous-%
groupe d'un groupe lisse (l'espace affine de dimension $r(N+1)$) , $v$ est
déterminé de manière unique (I.5.3) par la relation $\mathbb{V}\circ\mathbb{F}_{\mathbb{E}(M)|k} = p$ ;
nous allons exhiber un morphisme $\mathbb{V}$ satisfaisant à cette relation. On
montre que l'on a le diagramme commutatif suivant

\diagram{
\begin{tikzcd}[column sep=large]
\mathrm{Hom}_{\mathcal{D}}(M, \mathbb{W}_{n+1}(S)) \arrow[r, "\alpha \to \alpha\circ F"] \arrow[d, "\phi_S", "\wr"'] & \mathrm{Hom}_{\mathcal{D}}(M, \mathbb{W}_{N+1}(S,\sigma)) \arrow[d, "\phi_{S,\sigma}", "\wr"'] \\
\mathrm{Hom}_{k\text{-alg}}(A, \Gamma(\underline{\mathcal{O}}_S)) \arrow[r, "\mathbb{F}_{\mathbb{E}(M)|k}(S)"] & \mathrm{Hom}_{k\text{-alg}}(A \otimes_k (k,\sigma), \Gamma(\underline{\mathcal{O}}_S))
\end{tikzcd}
}

où $(S,\sigma)$ désigne le k-schéma $S \to \mathrm{Spec}(k) \overset{\sigma}{\to} \mathrm{Spec}(k)$ .

On construit alors un morphisme fonctoriel en $S$
\[ \mathbb{V} : \mathrm{Hom}_{\mathcal{D}}(M,\mathbb{W}_{N+1}(S,\sigma)) \to \mathrm{Hom}_{\mathcal{D}}(M,\mathbb{W}_{N+1}(S)) \]
\[ u \to u \circ V \ . \]

D'après ce qui a été dit, ce morphisme s'identifie à la Vers-%
chiebung. Comme $V$ est nilpotent sur $M$ , $\mathbb{V}$ est nilpotent et $\mathbb{E}(M)$
est unipotent.

\margin{\hand{52}}
